% =====================================================
% part8_multilinear.tex ― 第8部:多重線形代数
% =====================================================
\part{多重線形代数}

ここまでの主役は行列と、その分解であった。
最終理論部となる本部では、視点をさらに一段上げ、
\textbf{「線形」という概念そのものを解剖する}。
ベクトル空間から新しいベクトル空間を作る三つの普遍的な構成
――双対・商・テンソル積――を学び、最後に外積代数で、
本書第I部以来の旧友・行列式の「正体」を見届ける。
大学院入試で問われる線形代数の理論的話題は、おおむねここまでで覆われる。

% =====================================================
\chapter{双対空間 ― ベクトルを「測る」ものたちの空間}\label{chap:dual}
% =====================================================

\begin{goal}
\begin{itemize}
\item 線形形式(ベクトルを入れると数が出てくる線形な装置)の定義と実例を
知り、$\R^n$ 上の線形形式が\textbf{行ベクトル}と同一視できることを
説明できるようになる。
\item 与えられた基底の\textbf{双対基底}を計算できるようになり、
双対基底が $V^{*}$ の基底になることの証明を理解する。とくに、双対基底の
係数が\textbf{逆行列の行}から読み取れることを使えるようになる。
\item 線形写像 $f : V \to W$ から\textbf{双対写像} ${}^tf : W^{*} \to V^{*}$
が生まれ、その表現行列が\textbf{転置行列}になること
――第I部以来使い続けてきた転置の「正体」――を理解する。
\item ベクトルと線形形式の同一視が\textbf{基底の取り方に依存する}ことを、
反例で説明できるようになる。
\end{itemize}
\end{goal}

\section{動機 ― 「測る側」を主役にする}

本書は第I部の冒頭から、縦に書くベクトル(列)と横に書くベクトル(行)を
執拗に区別してきた。振り返れば、横書きのものはいつも
\textbf{「縦ベクトルを数に変える」役回り}で登場している。
連立一次方程式の $1$ 本の式
\[
a_1 x_1 + a_2 x_2 + \cdots + a_n x_n = b
\]
は「未知ベクトル $\vect{x}$ をこの係数の装置で測ったら $b$ になれ」
という条件だし、内積 $\langle \vect{a}, \vect{x} \rangle
= {}^t\vect{a}\,\vect{x}$(第V部)でも、測る側の $\vect{a}$ は
転置されて行になってから $\vect{x}$ に作用した。テストの合計点
(科目ごとの得点ベクトルに配点を掛けて足す)も、買い物の総額も、
みなこの形である。

本章では、この「測る装置」たちを一括して主役に据える。装置どうしは
足せるし定数倍もできるから、装置の全体は\textbf{それ自身が一つの
ベクトル空間}になる――これが双対空間である。そして章の後半では、
第V部・第VII部で大活躍した転置行列(正規方程式の ${}^t\!AA$、
直交行列の ${}^tQQ = E$、SVD の $U\Sigma\,{}^tV$)が、
\textbf{「測る側の世界で $A$ の代理として働く写像」の表現行列}である
という意味づけに決着をつける。

\section{線形形式と双対空間}

\begin{teigi}[線形形式・双対空間]
ベクトル空間 $V$ から $\R$ への線形写像
$f : V \to \R$ を $V$ 上の\textbf{線形形式}という。
線形形式の全体は、自然な和とスカラー倍
\[
(f + g)(\vect{v}) = f(\vect{v}) + g(\vect{v}), \qquad
(cf)(\vect{v}) = c\,f(\vect{v})
\]
によってそれ自身ベクトル空間になる。これを $V$ の\textbf{双対空間}といい、
$V^{*}$ と書く。
\end{teigi}

\begin{rei}[線形形式の実例]\label{rei:dual-examples}
(1) $\R^3$ 上の $f(\vect{x}) = 2x_1 - x_2 + 3x_3$ は線形形式である。
和とスカラー倍を保つことは、成分ごとの分配法則そのものである。\\
(2) 多項式の空間 $\R[x]_2$ 上で、「$x = 1$ での値を読む」装置
$\mathrm{ev}_1(p) = p(1)$ は線形形式である:
$\mathrm{ev}_1(p + q) = (p + q)(1) = p(1) + q(1)$。
「積分する」装置 $I(p) = \int_0^1 p(x)\,dx$ も、積分の線形性から
線形形式である。\\
(3) $2 \times 2$ 行列の空間 $M_2$ 上で、トレース
$\mathrm{tr}\,A = a_{11} + a_{22}$ は線形形式である。
どの例でも「入力はベクトル、出力は $1$ 個の数、そして一次で斉次」が
共通している。
\end{rei}

\begin{chui}
ノルム $\|\vect{v}\|$ は「ベクトルから数」だが線形形式では\textbf{ない}:
$\vect{v} \neq \vect{0}$ のとき $\|-\vect{v}\| = \|\vect{v}\|
\neq -\|\vect{v}\|$ となり、スカラー倍を保たない。
$q(\vect{x}) = x_1^2$ のような二次式も $q(2\vect{x}) = 4q(\vect{x})$ で
失格である。線形形式は「一次の斉次式」だけである。
\end{chui}

「ベクトルを入れると数が出てくる装置」たちが、また一つの
ベクトル空間をなす――この入れ子構造に最初は眩暈がするかもしれないが、
実例を見れば親しみやすい。$\R^n$ 上の線形形式は、必ず
\[
f(\vect{x}) = a_1 x_1 + a_2 x_2 + \cdots + a_n x_n
= \mat{a_1 & a_2 & \cdots & a_n} \vect{x}
\]
の形をしている。実際、$\vect{x} = x_1\vect{e}_1 + \cdots + x_n\vect{e}_n$
と標準基底で展開して線形性を使えば
\[
f(\vect{x}) = x_1 f(\vect{e}_1) + \cdots + x_n f(\vect{e}_n)
\]
となるから、$a_i = f(\vect{e}_i)$ と置けばよい。すなわち
\textbf{$\R^n$ の双対空間の要素とは、行ベクトルのことである}。
本書でずっと「縦ベクトル(列)」と「横ベクトル(行)」を区別してきたのは、
記法の趣味ではない。\textbf{列は測られるもの、行は測るもの}という、
役割の違う二種類の住人だったのである。

\section{双対基底}

$f(\vect{e}_i)$ たちが $f$ を決める、という上の観察は、
$\R^n$ に限らずどのベクトル空間でも通用する
(定理 \ref{thm:linmap-basis}:線形写像は基底の行き先で決まる)。
これを装置の側から言い直すと、「第 $i$ 座標だけを読み取る装置」たちが
$V^{*}$ の部品一式になる。

\begin{teigi}[双対基底]
$V$ の基底 $\vect{e}_1, \dots, \vect{e}_n$ に対し、
\[
f_i(\vect{e}_j) = \begin{cases} 1 & (i = j) \\ 0 & (i \neq j) \end{cases}
\]
で定まる線形形式 $f_1, \dots, f_n$ を、この基底の\textbf{双対基底}という
(基底の行き先を指定したので、各 $f_i$ は定理 \ref{thm:linmap-basis} により
一意に定まる)。
\end{teigi}

$f_i$ は「ベクトルの第 $i$ 座標を読み取る装置」である。
実際、$\vect{v} = c_1\vect{e}_1 + \cdots + c_n\vect{e}_n$ に対し、線形性から
$f_i(\vect{v}) = c_1 f_i(\vect{e}_1) + \cdots + c_n f_i(\vect{e}_n) = c_i$
となる。名前のとおり、双対基底は本当に $V^{*}$ の基底になる。

\begin{teiri}[双対基底は $V^{*}$ の基底である]\label{thm:dual-basis}
$f_1, \dots, f_n$ は $V^{*}$ の基底である。したがって
$\dim V^{*} = \dim V$。
\end{teiri}

\noindent\textbf{(証明)}\quad
\textbf{(線形独立)}\ $c_1 f_1 + \cdots + c_n f_n = o$(零形式)とする。
両辺を $\vect{e}_j$ で評価すると、左辺は $c_j$、右辺は $0$。
よってすべての $c_j = 0$。

\textbf{(生成)}\ 任意の $f \in V^{*}$ に対し、$c_i = f(\vect{e}_i)$ と置いて
$g = c_1 f_1 + \cdots + c_n f_n$ を作ると、
$g(\vect{e}_j) = c_j = f(\vect{e}_j)$ ですべての基底ベクトルの行き先が
一致するから、定理 \ref{thm:linmap-basis} より $f = g$。すなわち
\[
f = f(\vect{e}_1)\,f_1 + \cdots + f(\vect{e}_n)\,f_n
\]
と展開できる。展開係数は「$f$ に基底ベクトルを食わせた値」である。
$\blacksquare$

\begin{chui}
双対基底の $f_i$ は、相方の $\vect{e}_i$ \textbf{だけでは決まらない}。
「$\vect{e}_i$ 以外の基底ベクトルをすべて $0$ に送る」という条件が、
基底の\textbf{残りの全員}に依存するからである。たとえば $\R^2$ の基底
$\{\vect{e}_1, \vect{e}_2\}$ の双対基底では $f_1(x, y) = x$ だが、
基底 $\{\vect{e}_1,\ \vect{e}_1 + \vect{e}_2\}$ の双対基底では
$f_1(x, y) = x - y$ になる($f_1(\vect{e}_1) = 1$,
$f_1(\vect{e}_1 + \vect{e}_2) = 0$ を解けばよい)。
第 $2$ の基底ベクトルを取り替えただけで、$\vect{e}_1$ の「相方」まで
変わるのである。
\end{chui}

\begin{itembox}[l]{深掘り:価格ベクトルは双対空間の住人である}
双対空間の最も身近な実例は経済学にある。$n$ 種類の財の数量を並べた
\textbf{財ベクトル} $\vect{x} = \mat{x_1 \\ \vdots \\ x_n}$ は $\R^n$ の要素である。
一方、各財の価格 $p_1, \dots, p_n$ は、財ベクトルに\textbf{総額という数を
対応させる装置}
\[
\vect{x} \mapsto p_1 x_1 + \cdots + p_n x_n
\]
として働く。これはまさに線形形式であり、
\textbf{価格ベクトルは財空間そのものではなく、その双対空間に住んでいる}。
財と価格は次元(単位)が違う――個数と円/個――のだから、
同じ空間に住んでいないのはむしろ当然である。
物理学でも同様に、速度ベクトルに対する運動量、変位に対する力(仕事を
返す装置)は双対側の住人であり、量子力学のブラ・ケット記法
$\langle \phi | \psi \rangle$ は「双対の要素 $\langle \phi|$ で
ベクトル $|\psi\rangle$ を測る」ことを明示する記法である。
\textbf{「測るもの」と「測られるもの」を区別する文法}が双対性であり、
一度この眼鏡をかけると、応用分野の数式の意味が格段に読みやすくなる。
\end{itembox}

\begin{mondai}\label{q:dual}
$\R^2$ の基底 $\vect{v}_1 = \mat{1 \\ 1}$, $\vect{v}_2 = \mat{1 \\ -1}$ の
双対基底 $f_1, f_2$ を、$f_i(x, y) = ax + by$ の形で具体的に求めよ。
\end{mondai}

\section{双対写像 ― 転置行列の正体}

線形写像 $f : V \to W$ があると、$W$ 側の測定装置 $g \in W^{*}$ は、
「まず $f$ で送ってから測る」ことで $V$ 側の装置に変わる。
合成 $g \circ f : V \to \R$ は線形写像の合成だから線形であり
(定理 \ref{thm:linmap-comp})、確かに $V$ 上の線形形式である。

\begin{teigi}[双対写像]\label{def:dual-map}
線形写像 $f : V \to W$ に対し、
\[
({}^tf)(g) = g \circ f \qquad (g \in W^{*})
\]
で定まる写像 ${}^tf : W^{*} \to V^{*}$ を $f$ の\textbf{双対写像}という。
${}^tf$ 自身も線形写像である。実際、任意の $\vect{v} \in V$ で
\[
\bigl((g + h) \circ f\bigr)(\vect{v}) = g(f(\vect{v})) + h(f(\vect{v})), \qquad
\bigl((cg) \circ f\bigr)(\vect{v}) = c\,g(f(\vect{v}))
\]
だから、${}^tf(g + h) = {}^tf(g) + {}^tf(h)$、${}^tf(cg) = c\,{}^tf(g)$。
\end{teigi}

向きに注意してほしい。$f$ は $V$ から $W$ へ進むが、双対写像は
$W^{*}$ から $V^{*}$ へと\textbf{逆向き}に進む。測る装置は、
写像をさかのぼって運ばれるのである。記号 ${}^tf$ を使った理由は次の定理にある。

\begin{teiri}[双対写像の表現行列は転置行列]\label{thm:dual-transpose}
$f : \R^n \to \R^m$ の(標準基底に関する)表現行列を $A$ とすると、
双対基底に関する ${}^tf : (\R^m)^{*} \to (\R^n)^{*}$ の表現行列は
${}^t\!A$ である。
\end{teiri}

\noindent\textbf{(証明)}\quad
$g \in (\R^m)^{*}$ を行ベクトル $\vect{w} = \mat{w_1 & \cdots & w_m}$ と
同一視する($g(\vect{y}) = \vect{w}\vect{y}$)。合成を計算すると
\[
\bigl(({}^tf)(g)\bigr)(\vect{x}) = g(A\vect{x}) = \vect{w}(A\vect{x})
= (\vect{w}A)\vect{x}
\]
だから、行ベクトルの言葉では ${}^tf$ は「\textbf{右から $A$ を掛ける}」
操作 $\vect{w} \mapsto \vect{w}A$ である。双対基底に関する座標は、
行ベクトルの成分を縦に並べた列だから、座標の変換は
\[
{}^t(\vect{w}A) = {}^t\!A\,{}^t\vect{w},
\]
すなわち「座標の列に ${}^t\!A$ を掛ける」ことになる。表現行列は座標に
掛かる行列だった(定理 \ref{thm:repmat-coord})から、
${}^tf$ の表現行列は ${}^t\!A$ である。 $\blacksquare$

これが\textbf{転置行列の正体}である:${}^t\!A$ とは、
「測る側の世界で $A$ の代理として働く写像」の行列にほかならない。
この眼鏡をかけると、本書でくり返し使った等式
\[
\langle A\vect{x}, \vect{y} \rangle = \langle \vect{x}, {}^t\!A\vect{y} \rangle
\quad (\text{第V部})
\]
は「$\vect{x}$ を $A$ で送ってから測っても、測る装置を ${}^t\!A$ で
引き戻してから測っても、結果は同じ」という、双対写像の定義の言い換え
そのものになる。また第II部の定理 \ref{thm:rank-transpose}
($\mathrm{rank}\,{}^t\!A = \mathrm{rank}\,A$)は、
「写像 $f$ とその双対 ${}^tf$ は同じ階数をもつ」と読み直せる。

\section{双対の双対 ― 元の世界に戻る}

双対の双対 $V^{**}$、すなわち「測るものを測るもの」の空間はどうなるか。
$\vect{v} \in V$ に対して「$f$ を $f(\vect{v})$ にうつす装置」
$\iota(\vect{v}) : V^{*} \to \R$ を対応させると、
$\iota(\vect{v})$ は $V^{*}$ 上の線形形式(定義の演算がそう作ってある)
であり、対応 $\iota : V \to V^{**}$ は線形写像になる。
しかも $\iota$ は\textbf{単射}である:$\vect{v} \neq \vect{0}$ なら
$f(\vect{v}) \neq 0$ となる線形形式 $f$ が存在する
(例題 \ref{rei:dual-sep} の分離補題)から、
$\iota(\vect{v})$ は零形式でない。定理 \ref{thm:dual-basis} を二回使えば
$\dim V^{**} = \dim V^{*} = \dim V$ なので、単射線形写像 $\iota$ は
次元定理(定理 \ref{thm:dimension})により全射でもあり、同型になる。
「測るものを測るもの」は元の世界に戻ってくるのである。

重要なのは、この同一視 $\iota$ が\textbf{基底を一切選ばずに}作れたこと
である。$V$ と $V^{*}$ も次元が等しいから同型ではあるが、そちらの
同一視は基底の取り方に依存する(次節の反例)。
$V \cong V^{**}$ だけが、誰が作っても同じになる「自然な」同型である。

\section{反例 ― ありがちな誤解を正す}

\begin{hanrei}[ベクトルと線形形式の同一視は基底に依存する]\label{hanrei:dual-ident}
「$\R^n$ の双対は行ベクトル、つまり成分は同じで向きが違うだけ。
それなら $V$ と $V^{*}$ は同じものでは?」――この同一視は、
\textbf{どの基底で成分を読むか}に依存する。基底 $\{\vect{b}_i\}$ を選び、
「座標が同じもの同士」$\sum c_i \vect{b}_i \mapsto \sum c_i g_i$
($\{g_i\}$ は双対基底)で対応させてみよう。
$\R^2$ のベクトル $\vect{v} = \mat{2 \\ 0}$ に対し、
標準基底 $\{\vect{e}_1, \vect{e}_2\}$ を使うと
$\vect{v} = 2\vect{e}_1$ の相手は $2f_1$、すなわち $f(x, y) = 2x$ である。
ところが基底 $\vect{v}_1 = \mat{2 \\ 0}$, $\vect{v}_2 = \mat{0 \\ 1}$ を
使うと、$\vect{v} = 1 \cdot \vect{v}_1$ の相手は双対基底の $g_1$、
つまり $g_1(x, y) = \dfrac{x}{2}$ である($g_1(\vect{v}_1) = 1$ に注意)。
$2x$ と $\dfrac{x}{2}$ ――同じベクトルの「相手」が、基底の選び方で
$4$ 倍も違ってしまった。基底ベクトルを $2$ 倍すると座標は
$\dfrac{1}{2}$ 倍、双対基底も $\dfrac{1}{2}$ 倍になる(逆向きに変わる)
ことが原因である。単位の取り替え(メートル→センチ)で
財の数量と価格が逆方向に換算されるのと同じ現象であり、
「列と行はあくまで別世界の住人」と区別しておくのが安全である。
なお第V部のように\textbf{内積を一つ固定}すれば、
$\vect{a} \mapsto \langle \vect{a}, \cdot \rangle$ という基底によらない
同一視が作れる。行と列を無造作に行き来できたのは、
暗黙に標準内積を使っていたからである。
\end{hanrei}

\begin{hanrei}[線形従属な組には値を自由に指定できない]\label{hanrei:dual-dep}
双対基底の構成(基底の行き先を自由に指定する)は、
\textbf{基底でない組には通用しない}。たとえば
$\vect{v}_1 = \mat{1 \\ 2}$, $\vect{v}_2 = \mat{2 \\ 4}$ に対し、
$f(\vect{v}_1) = 1$, $f(\vect{v}_2) = 3$ となる線形形式は存在しない。
$\vect{v}_2 = 2\vect{v}_1$ だから、線形性より
$f(\vect{v}_2) = 2f(\vect{v}_1) = 2 \neq 3$ と、値が勝手に決まって
しまうからである。「基底の行き先は自由、従属な組の行き先は不自由」は
線形写像一般の鉄則であり、定理 \ref{thm:linmap-basis} が
「基底」を仮定していたのはこのためである。双対基底の構成が
基底を要求する理由もここにある。
\end{hanrei}

\section{例題}

双対基底の計算・線形形式の展開・双対写像を練習する。
以下は\textbf{解答つきの例題}である。手を動かして追ってほしい。

\begin{rei}[例題・基本(行ベクトル表示)]\label{rei:dual-calc}
$\R^3$ 上の線形形式 $f(x, y, z) = 2x - y + 3z$ を行ベクトルで表し、
$\vect{u} = {}^t(1, 1, 1)$, $\vect{v} = {}^t(2, 0, -1)$ での値と、
$f(\vect{u} + \vect{v})$ を求めよ。\\
\textbf{解.}\ $f$ の行ベクトルは $\mat{2 & -1 & 3}$ である。
\[
f(\vect{u}) = 2 - 1 + 3 = 4, \qquad f(\vect{v}) = 4 + 0 - 3 = 1.
\]
$\vect{u} + \vect{v} = {}^t(3, 1, 0)$ だから
$f(\vect{u} + \vect{v}) = 6 - 1 + 0 = 5 = f(\vect{u}) + f(\vect{v})$。
線形性が数字で確かめられた。
\end{rei}

\begin{rei}[例題・基本(双対基底の計算)]\label{rei:dual-basis2}
$\R^2$ の基底 $\vect{v}_1 = \mat{1 \\ 0}$, $\vect{v}_2 = \mat{1 \\ 1}$ の
双対基底 $g_1, g_2$ を求め、$\vect{w} = \mat{3 \\ 5}$ を
$\vect{w} = g_1(\vect{w})\vect{v}_1 + g_2(\vect{w})\vect{v}_2$ と
展開してみよ。\\
\textbf{解.}\ $g_1(x, y) = ax + by$ に $g_1(\vect{v}_1) = 1$,
$g_1(\vect{v}_2) = 0$ を課すと $a = 1$, $a + b = 0$ より
$g_1(x, y) = x - y$。同様に $g_2$ は $a = 0$, $a + b = 1$ より
$g_2(x, y) = y$。展開係数は
$g_1(\vect{w}) = 3 - 5 = -2$, $g_2(\vect{w}) = 5$ で、検算すると
\[
-2\mat{1 \\ 0} + 5\mat{1 \\ 1} = \mat{3 \\ 5} = \vect{w}
\]
となり、確かに座標を読み取れている。
\end{rei}

\begin{rei}[例題・基本(多項式の空間で)]\label{rei:dual-poly}
$V = \R[x]_1$($1$ 次以下の多項式)上の線形形式
$\mathrm{ev}_0(p) = p(0)$, $\mathrm{ev}_1(p) = p(1)$ を考える。
$\{\mathrm{ev}_0, \mathrm{ev}_1\}$ が双対基底になるような $V$ の基底
$\{p_1, p_2\}$ を求めよ。\\
\textbf{解.}\ 条件は $p_1(0) = 1$, $p_1(1) = 0$ と
$p_2(0) = 0$, $p_2(1) = 1$。$1$ 次式でこれを満たすのは
\[
p_1(x) = 1 - x, \qquad p_2(x) = x
\]
である。このとき任意の $p \in V$ は
\[
p = \mathrm{ev}_0(p)\,p_1 + \mathrm{ev}_1(p)\,p_2
= p(0)(1 - x) + p(1)\,x
\]
と展開される――「両端の値から直線を復元する」補間公式であり、
双対基底の展開公式(定理 \ref{thm:dual-basis} の証明)の実例になっている。
\end{rei}

\begin{rei}[例題・標準(双対基底は逆行列の行)]\label{rei:dual-inv}
$\R^3$ の基底 $\vect{v}_1 = \mat{1 \\ 0 \\ 0}$,
$\vect{v}_2 = \mat{1 \\ 1 \\ 0}$, $\vect{v}_3 = \mat{1 \\ 1 \\ 1}$ の
双対基底を求めよ。\\
\textbf{解.}\ 双対基底 $g_i$ の行ベクトルを縦に積んだ行列を $F$、
$\vect{v}_j$ を並べた行列を $P = \mat{1 & 1 & 1 \\ 0 & 1 & 1 \\ 0 & 0 & 1}$
とすると、$(FP)_{ij} = g_i(\vect{v}_j)$ だから、双対基底の条件は
$FP = E$、すなわち\textbf{$F = P^{-1}$(双対基底は逆行列の行)}である。
掃き出しで
\[
P^{-1} = \mat{1 & -1 & 0 \\ 0 & 1 & -1 \\ 0 & 0 & 1}
\]
だから、
\[
g_1(x, y, z) = x - y, \qquad g_2(x, y, z) = y - z, \qquad
g_3(x, y, z) = z.
\]
検算:$g_1(\vect{v}_1) = 1$, $g_1(\vect{v}_2) = 1 - 1 = 0$,
$g_1(\vect{v}_3) = 0$ など、すべて条件を満たす。
第I部の逆行列計算が、そのまま双対基底の計算装置になる。
\end{rei}

\begin{rei}[例題・標準(双対写像の計算)]\label{rei:dual-map}
$f : \R^2 \to \R^2$ の表現行列を $A = \mat{1 & 2 \\ 3 & 4}$、
$g(x, y) = x + y$ とする。$({}^tf)(g) = g \circ f$ を直接計算し、
その係数が ${}^t\!A$ を掛けて得られることを確かめよ。\\
\textbf{解.}\ 直接計算では
\[
(g \circ f)(x, y) = g(x + 2y,\ 3x + 4y) = (x + 2y) + (3x + 4y) = 4x + 6y.
\]
一方、$g$ の係数の列は ${}^t(1, 1)$ で、
\[
{}^t\!A\mat{1 \\ 1} = \mat{1 & 3 \\ 2 & 4}\mat{1 \\ 1} = \mat{4 \\ 6}
\]
と、確かに $4x + 6y$ の係数が得られる(定理 \ref{thm:dual-transpose})。
行ベクトルで書けば $\mat{1 & 1}A = \mat{4 & 6}$:
「右から $A$ を掛ける」操作である。
\end{rei}

\begin{rei}[例題・標準(証明:分離補題)]\label{rei:dual-sep}
$\vect{v} \in V$ について、「すべての $f \in V^{*}$ で $f(\vect{v}) = 0$」
ならば $\vect{v} = \vect{0}$ であることを示せ。\\
\textbf{解.}\ 対偶を示す。$\vect{v} \neq \vect{0}$ とすると、
$1$ 本だけの組 $\{\vect{v}\}$ は線形独立だから、基底
$\vect{v}, \vect{u}_2, \dots, \vect{u}_n$ に延長できる
(定理 \ref{thm:dim-extend})。その双対基底の第 $1$ 要素 $f_1$ は
$f_1(\vect{v}) = 1 \neq 0$ を満たす。すなわち $\vect{v} \neq \vect{0}$
なら、それを検出する測定装置が必ず存在する。
(この補題が、前節の $\iota : V \to V^{**}$ の単射性の根拠である。)
\end{rei}

\begin{matome}
\begin{itemize}
\item $V$ 上の\textbf{線形形式}(ベクトルを数にうつす線形写像)の全体が
\textbf{双対空間} $V^{*}$。$\R^n$ の線形形式は\textbf{行ベクトル}
$\mat{a_1 & \cdots & a_n}$($a_i = f(\vect{e}_i)$)と同一視できる。
\item 基底 $\vect{e}_1, \dots, \vect{e}_n$ の\textbf{双対基底}
$f_1, \dots, f_n$($f_i(\vect{e}_j) = \delta$ 型の条件)は $V^{*}$ の
基底になり(定理 \ref{thm:dual-basis})、$\dim V^{*} = \dim V$。
$f_i$ は第 $i$ 座標の読み取り装置で、
$f = \sum f(\vect{e}_i) f_i$ と展開できる。
$\R^n$ では\textbf{双対基底は $P = (\vect{v}_1\ \cdots\ \vect{v}_n)$ の
逆行列の行}(例題 \ref{rei:dual-inv})。
\item 線形写像 $f : V \to W$ の\textbf{双対写像}
${}^tf : W^{*} \to V^{*}$($g \mapsto g \circ f$)の表現行列は
\textbf{転置行列} ${}^t\!A$(定理 \ref{thm:dual-transpose})。
$\langle A\vect{x}, \vect{y} \rangle = \langle \vect{x},
{}^t\!A\vect{y} \rangle$ はその言い換えである。
\item $V \cong V^{**}$ は基底によらず自然に成り立つが、
$V \cong V^{*}$ の同一視は\textbf{基底(または内積)の選択に依存する}
(反例 \ref{hanrei:dual-ident})。
\item 双対基底の構成には「基底であること」が必須:従属な組には
値を自由に指定できない(反例 \ref{hanrei:dual-dep})。
\end{itemize}
\end{matome}

\section*{章末問題}

\begin{mondai}[基本]\label{q:dual-ex1}
$\R^2$ 上の線形形式 $f(x, y) = 3x - 2y$ を行ベクトルで表せ。また
$\vect{u} = {}^t(1, 2)$, $\vect{v} = {}^t(-1, 1)$ について
$f(\vect{u})$, $f(\vect{v})$, $f(\vect{u} + \vect{v})$ を求め、
線形性を数字で確かめよ。
\end{mondai}

\begin{mondai}[基本]\label{q:dual-ex2}
$\R^2$ の基底 $\vect{v}_1 = \mat{2 \\ 1}$, $\vect{v}_2 = \mat{1 \\ 1}$ の
双対基底 $g_1, g_2$ を $g_i(x, y) = ax + by$ の形で求めよ。
\end{mondai}

\begin{mondai}[基本]\label{q:dual-ex3}
$\R[x]_1$ 上の線形形式 $I(p) = \displaystyle\int_0^1 p(x)\,dx$ を、
例題 \ref{rei:dual-poly} の $\mathrm{ev}_0, \mathrm{ev}_1$ の線形結合で
表せ(結果を台形公式と見比べよ)。
\end{mondai}

\begin{mondai}[基本]\label{q:dual-ex4}
$2 \times 2$ 行列の空間 $M_2$ 上のトレース $\mathrm{tr}$ が線形形式で
あることを確かめ、基底 $E_{11}, E_{12}, E_{21}, E_{22}$
($E_{ij}$ は $(i,j)$ 成分だけが $1$ の行列)に関する係数
$\mathrm{tr}(E_{ij})$ をすべて求めよ。また
$\mathrm{tr}\mat{3 & 5 \\ 2 & 7}$ を計算せよ。
\end{mondai}

\begin{mondai}[標準]\label{q:dual-ex5}
$\R^3$ の基底 $\vect{v}_1 = \mat{1 \\ 1 \\ 0}$,
$\vect{v}_2 = \mat{0 \\ 1 \\ 1}$, $\vect{v}_3 = \mat{1 \\ 0 \\ 1}$ の
双対基底を、例題 \ref{rei:dual-inv} の方法(逆行列の行)で求めよ。
\end{mondai}

\begin{mondai}[標準]\label{q:dual-ex6}
$f : \R^2 \to \R^2$ の表現行列を $A = \mat{1 & 1 \\ 0 & 1}$、
$g(x, y) = 2x + 3y$ とする。$({}^tf)(g)$ を合成の直接計算と
${}^t\!A$ の掛け算の二通りで求め、一致を確かめよ。
\end{mondai}

\begin{mondai}[標準]\label{q:dual-ex7}
$V = \R[x]_2$ 上の三つの線形形式
$\mathrm{ev}_{-1}(p) = p(-1)$, $\mathrm{ev}_0(p) = p(0)$,
$\mathrm{ev}_1(p) = p(1)$ に対し、
$\{\mathrm{ev}_{-1}, \mathrm{ev}_0, \mathrm{ev}_1\}$ が双対基底になる
$V$ の基底 $\{p_{-1}, p_0, p_1\}$ を求めよ
(ヒント:$p_{-1}$ は $x = 0, 1$ で零になる $2$ 次式)。
\end{mondai}

\begin{mondai}[標準]\label{q:dual-ex8}
$\R^2$ 上の線形形式 $f$ が $f\mat{1 \\ 1} = 3$,
$f\mat{1 \\ -1} = 1$ を満たすとき、$f(x, y)$ を求めよ
(ヒント:問 \ref{q:dual} の双対基底で $f$ を展開する)。
\end{mondai}

\begin{mondai}[発展]\label{q:dual-ex9}
部分空間 $W \subset V$ に対し、$W$ 上で恒等的に零になる線形形式の全体
\[
W^{0} = \{\, f \in V^{*} \mid f(\vect{w}) = 0 \ (\forall \vect{w} \in W) \,\}
\]
を $W$ の\textbf{零化空間}という。
$V = \R^3$, $W = \langle {}^t(1, 1, 0),\ {}^t(0, 1, 1) \rangle$ について
$W^{0}$ を具体的に求め、$\dim W^{0} = \dim V - \dim W$ が成り立つことを
確かめよ。
\end{mondai}

\begin{mondai}[証明]\label{q:dual-ex10}
$\vect{v}_1, \dots, \vect{v}_k \in V$ について、次の同値を示せ:
$f_i(\vect{v}_j) = \begin{cases} 1 & (i = j) \\ 0 & (i \neq j)\end{cases}$
を満たす $f_1, \dots, f_k \in V^{*}$ が存在する
$\iff$ $\vect{v}_1, \dots, \vect{v}_k$ は線形独立である
(ヒント:$\Leftarrow$ は基底への延長(定理 \ref{thm:dim-extend})と
双対基底、$\Rightarrow$ は関係式 $\sum c_j \vect{v}_j = \vect{0}$ に
$f_i$ を施す)。
\end{mondai}

\begin{mondai}[証明]\label{q:dual-ex11}
$V$ を $n$ 次元ベクトル空間、$f \in V^{*}$ を零形式でない線形形式とする。
$\dim \mathrm{Ker}\,f = n - 1$ を示せ
(ヒント:次元定理(定理 \ref{thm:dimension})。$\mathrm{Im}\,f$ は
$\R$ の部分空間である)。
\end{mondai}

\bigskip
{\small
\noindent\textbf{【章末問題の方針】}\quad
まず自力で取り組み、行き詰まったら下の方針を見よ。記述による完全解答は
巻末「問の解答」にある。
\begin{itemize}
\item \textbf{問 \ref{q:dual-ex1}}:行ベクトルは係数を並べるだけ。
値は代入して計算する。
\item \textbf{問 \ref{q:dual-ex2}}:$g_i(\vect{v}_j)$ の $4$ つの条件から
連立方程式。または $P = \mat{2 & 1 \\ 1 & 1}$ の逆行列の行
(例題 \ref{rei:dual-inv})。
\item \textbf{問 \ref{q:dual-ex3}}:$I = I(p_1)\mathrm{ev}_0
+ I(p_2)\mathrm{ev}_1$(定理 \ref{thm:dual-basis} の展開公式を
$V^{*}$ の基底 $\{\mathrm{ev}_0, \mathrm{ev}_1\}$ に適用)。
$p_1 = 1 - x$, $p_2 = x$ を積分する。
\item \textbf{問 \ref{q:dual-ex4}}:トレースの線形性は成分の計算。
係数は対角の $E_{ij}$ だけが $1$。
\item \textbf{問 \ref{q:dual-ex5}}:$P = (\vect{v}_1\ \vect{v}_2\ \vect{v}_3)$
の逆行列を掃き出しで求め、その行を読む。
\item \textbf{問 \ref{q:dual-ex6}}:例題 \ref{rei:dual-map} と同じ手順。
行ベクトル $\mat{2 & 3}$ に右から $A$ を掛けてもよい。
\item \textbf{問 \ref{q:dual-ex7}}:各 $p_i$ は他の $2$ 点で零になる
$2$ 次式を作り、残る $1$ 点で値 $1$ に正規化する
(ラグランジュ補間の基本式)。
\item \textbf{問 \ref{q:dual-ex8}}:$f = 3f_1 + 1f_2$
($f_1, f_2$ は問 \ref{q:dual} の答え)。
\item \textbf{問 \ref{q:dual-ex9}}:$f = ax + by + cz$ に
$2$ つの条件を課して解く。自由度の個数が次元になる。
\item \textbf{問 \ref{q:dual-ex10}}:両方向とも本文の道具
(双対基底・線形性)で書ける。
\item \textbf{問 \ref{q:dual-ex11}}:$f \neq o$ なら
$\mathrm{Im}\,f \neq \{0\}$。$\R$ の部分空間は $\{0\}$ と $\R$ しかない。
\end{itemize}
}

% =====================================================
\chapter{商空間 ― 「違いを無視する」という構成}\label{chap:quotient}
% =====================================================

\section{部分空間で割る}

整数を $3$ で割った余りで分類すると、無限にある整数が
$\{0, 1, 2\}$ の三クラスに「圧縮」され、クラスどうしで
矛盾なく足し算ができた(合同式)。同じことをベクトル空間で行う。

\begin{teigi}[商空間]
ベクトル空間 $V$ とその部分空間 $W$ に対し、
「差が $W$ に属するベクトルどうしを同一視する」と定める:
\[
\vect{u} \sim \vect{v} \iff \vect{u} - \vect{v} \in W.
\]
$\vect{v}$ と同一視されるベクトル全体 $\vect{v} + W$ を
($\vect{v}$ の属する)\textbf{類}といい、類全体の集合を $V/W$ と書く。
類どうしの演算を
\[
(\vect{u} + W) + (\vect{v} + W) = (\vect{u} + \vect{v}) + W, \qquad
c(\vect{v} + W) = c\vect{v} + W
\]
で定めると、$V/W$ はベクトル空間になる(代表の取り方によらず
矛盾なく定まることが要点である)。これを\textbf{商空間}という。
\end{teigi}

幾何的なイメージ:$V = \R^2$、$W$ を原点を通る直線とすると、
類 $\vect{v} + W$ は「$W$ に平行な直線」である。商空間 $\R^2/W$ とは、
\textbf{平行線の一本一本を「一つの点」とみなした世界}であり、
直線の束は $1$ 個のパラメータで番号付けできるから $1$ 次元になる。
一般に次が成り立つ。

\begin{teiri}
$V$ が有限次元のとき、$\dim(V/W) = \dim V - \dim W$。
\end{teiri}

\section{準同型定理 ― 次元定理の完成形}

\begin{teiri}[準同型定理]
線形写像 $f : V \to U$ に対して、
\[
V/\mathrm{Ker}\,f \;\cong\; \mathrm{Im}\,f.
\]
同型写像は $\vect{v} + \mathrm{Ker}\,f \mapsto f(\vect{v})$ で与えられる。
\end{teiri}

読み方:$f$ で同じ行き先につぶれるベクトルたち(差が核に属するもの)を
ひとまとめにすれば、$f$ は\textbf{一対一}になる。
「つぶれる分をあらかじめ割っておけば、写像は同型になる」のである。
両辺の次元をとれば
\[
\dim V - \dim \mathrm{Ker}\,f = \dim \mathrm{Im}\,f
\]
となり、第III部の次元定理が再現される。次元定理は準同型定理の「影」
(次元だけ見た姿)だったのである。

\begin{itembox}[l]{深掘り:割って同型を作る ― 代数学の普遍文法}
「写像 $f$ があれば、定義域を $\mathrm{Ker}\,f$ で割ったものが
$\mathrm{Im}\,f$ と同型になる」という形の定理は、ベクトル空間に限らず、
群でも環でも(将来代数学で学ぶあらゆる構造で)成り立つ。
合同式 $\bmod\ 3$ は「整数を、$3$ の倍数という核で割る」操作であり、
角度が $360^\circ$ で一周するのは「実数を $360$ の倍数で割った商」に
住んでいるからである。\textbf{構造を保つ写像(準同型)・核・商・同型}という
四点セットは、現代代数学の文法そのものであり、線形代数はその文法を
最も計算しやすい場所で練習させてくれる。商空間はまた、
「興味のない自由度を消して本質だけ残す」操作として、
幾何学(射影空間)や物理(ゲージ理論)でも繰り返し現れる。
\end{itembox}

\begin{mondai}\label{q:quot}
線形写像 $f : \R^3 \to \R^2$ を $f(x, y, z) = \mat{x - y \\ y - z}$ で定める。
$\mathrm{Ker}\,f$ を求め、$f$ が全射であることを示し、
準同型定理の両辺の次元が一致することを確かめよ。
\end{mondai}

% =====================================================
\chapter{テンソル積 ― 空間の「掛け算」}\label{chap:tensor}
% =====================================================

\section{動機 ― 双線形性をほどく}

二つの変数のそれぞれについて線形な写像(\textbf{双線形写像})は、
本書に何度も登場した:内積 $\langle \vect{u}, \vect{v} \rangle$、
行列の積、二次形式。双線形写像は線形写像ではない
($\langle 2\vect{u}, 2\vect{v} \rangle = 4\langle \vect{u}, \vect{v} \rangle$)ので、
これまでの理論がそのままでは使えない。そこで、
\textbf{双線形を線形に変換する万能の装置}を作る。それがテンソル積である。

\begin{teigi}[テンソル積]
ベクトル空間 $V$($\dim = m$、基底 $\vect{e}_1, \dots, \vect{e}_m$)と
$W$($\dim = n$、基底 $\vect{f}_1, \dots, \vect{f}_n$)に対し、
$mn$ 個の形式的な記号 $\vect{e}_i \otimes \vect{f}_j$ を基底とする
$mn$ 次元ベクトル空間を $V \otimes W$ と書き、\textbf{テンソル積空間}という。
さらに $\vect{u} = \sum_i u_i \vect{e}_i$, $\vect{v} = \sum_j v_j \vect{f}_j$
に対して
\[
\vect{u} \otimes \vect{v} = \sum_{i, j} u_i v_j\,(\vect{e}_i \otimes \vect{f}_j)
\]
と定める。記号 $\otimes$ は各変数について線形になる:
$(\vect{u}_1 + \vect{u}_2) \otimes \vect{v} = \vect{u}_1 \otimes \vect{v}
+ \vect{u}_2 \otimes \vect{v}$、
$(c\vect{u}) \otimes \vect{v} = c(\vect{u} \otimes \vect{v})
= \vect{u} \otimes (c\vect{v})$。
\end{teigi}

要点は二つ。第一に\textbf{次元は和でなく積}になる
($\dim(V \otimes W) = mn$)。第二に、$V \times W$ からの双線形写像と、
$V \otimes W$ からの\textbf{線形}写像が一対一に対応する
(双線形写像 $B$ に対し、$\vect{e}_i \otimes \vect{f}_j \mapsto
B(\vect{e}_i, \vect{f}_j)$ で線形写像が決まる)。
双線形の問題はテンソル積上の線形の問題に翻訳でき、
本書の全理論が適用可能になる――これがテンソル積の存在意義である。

\section{$\vect{u} \otimes \vect{v}$ の形に書けないテンソル}

$V \otimes W$ の要素で $\vect{u} \otimes \vect{v}$ の形のものを
\textbf{分解可能}という。重要なのは、
\textbf{一般のテンソルは分解可能とは限らない}ことである。
$\vect{u} \otimes \vect{v}$ の係数 $u_i v_j$ を並べた行列は
$\vect{u}\,{}^t\vect{v}$ という\textbf{階数 $1$} の行列になるから、
係数行列の階数が $2$ 以上のテンソルは分解できない
(第VII部の低ランク近似で活躍した「階数 $1$ の部品」が、
テンソルの言葉では「分解可能テンソル」だったのである)。

\begin{itembox}[l]{深掘り:量子もつれ ― 分解できないテンソルの物理}
量子力学では、系の状態はベクトル(状態ベクトル)で表され、
\textbf{二つの系を合わせた全体の状態空間は、各系の状態空間の
テンソル積}になる、というのが基本原理である。
次元が積で増える($2$ 状態系を $k$ 個合わせると $2^k$ 次元)ことが、
量子コンピュータの潜在能力の源泉である。

合成系の状態が $\vect{u} \otimes \vect{v}$ と分解可能なら、
「系 $1$ は状態 $\vect{u}$、系 $2$ は状態 $\vect{v}$」と
別々に記述できる。しかし上で見たとおり、
\textbf{分解できない状態}が存在する。たとえば
\[
\vect{e}_1 \otimes \vect{f}_1 + \vect{e}_2 \otimes \vect{f}_2
\]
(係数行列が単位行列で階数 $2$)。このとき二つの系は、
どちらか一方だけの状態を語ることができない一体の存在になっている。
これが\textbf{量子もつれ(エンタングルメント)}であり、
アインシュタインが「不気味な遠隔作用」と呼んで疑い、
後に実験で実証された(2022年ノーベル物理学賞の主題)現象の数学的正体である。
\textbf{線形代数の語彙(階数!)が、そのまま物理の最深部の語彙になっている}。
\end{itembox}

なお、行列の世界でのテンソル積は\textbf{クロネッカー積}
($A \otimes B$:$A$ の各成分 $a_{ij}$ を $a_{ij}B$ というブロックに
置き換えた行列)として実装され、$\dim$ が積になることが
ブロックの個数として目に見える。

\begin{mondai}\label{q:tensor}
$\R^2 \otimes \R^2$(基底 $\vect{e}_i \otimes \vect{e}_j$、$i, j = 1, 2$)
において、テンソル
$t = \vect{e}_1 \otimes \vect{e}_1 + \vect{e}_2 \otimes \vect{e}_2$ が
分解可能でない(どんな $\vect{u}, \vect{v}$ でも
$t = \vect{u} \otimes \vect{v}$ と書けない)ことを、係数を比較して示せ。
\end{mondai}

% =====================================================
\chapter{外積代数 ― 行列式の正体}\label{chap:wedge}
% =====================================================

\section{交代性を組み込んだ積}

テンソル積 $\vect{u} \otimes \vect{v}$ に
\textbf{交代性}(入れ替えると符号が反転する)を要求した積を考える。

\begin{teigi}[外積(ウェッジ積)]
ベクトル $\vect{u}, \vect{v} \in V$ に対し、双線形かつ
\[
\vect{u} \wedge \vect{v} = -\,\vect{v} \wedge \vect{u}
\qquad (\text{特に } \vect{u} \wedge \vect{u} = \vect{0})
\]
を満たす積 $\wedge$ を\textbf{外積}という。$k$ 本のウェッジ
$\vect{u}_1 \wedge \cdots \wedge \vect{u}_k$ の張る空間を
$\Lambda^k V$ と書く。$\Lambda^k V$ の基底は
$\vect{e}_{i_1} \wedge \cdots \wedge \vect{e}_{i_k}$($i_1 < \cdots < i_k$)
であり、次元は二項係数 $\binom{n}{k}$ になる
(同じ添字は消え、順序の入れ替えは符号に吸収されるため、
\textbf{添字の選び方}だけが自由度として残る)。
\end{teigi}

計算してみよう。$\R^2$ で $\vect{u} = \mat{a \\ c}$,
$\vect{v} = \mat{b \\ d}$ とすると、双線形性と交代性により
\[
\vect{u} \wedge \vect{v}
= (a\vect{e}_1 + c\vect{e}_2) \wedge (b\vect{e}_1 + d\vect{e}_2)
= ad\,(\vect{e}_1 \wedge \vect{e}_2) + cb\,(\vect{e}_2 \wedge \vect{e}_1)
= (ad - bc)\,\vect{e}_1 \wedge \vect{e}_2.
\]
\textbf{行列式が、計算せずとも勝手に現れた}。
$\vect{e}_1 \wedge \vect{e}_1$ と $\vect{e}_2 \wedge \vect{e}_2$ が消え、
$\vect{e}_2 \wedge \vect{e}_1 = -\vect{e}_1 \wedge \vect{e}_2$ と
符号が付く――第II部で行列式の定義に込めた「列の重複は許さない、
入れ替えは符号」という規則が、\textbf{積の代数規則として}実装されているからである。

\section{行列式の最終的な定義}

一般に $\dim \Lambda^n V = \binom{n}{n} = 1$、すなわち
\textbf{最高次の外積空間は $1$ 次元}である。
線形変換 $f : V \to V$ は $\Lambda^n V$ 上に
$\vect{u}_1 \wedge \cdots \wedge \vect{u}_n \mapsto
f(\vect{u}_1) \wedge \cdots \wedge f(\vect{u}_n)$ という線形変換を誘導するが、
$1$ 次元空間の線形変換は\textbf{ある定数倍}でしかありえない。

\begin{teiri}[行列式の基底によらない定義]
線形変換 $f$ が $\Lambda^n V$ 上に誘導する「定数倍」の定数が、
$f$ の行列式である:
\[
f(\vect{u}_1) \wedge \cdots \wedge f(\vect{u}_n)
= (\det f)\; \vect{u}_1 \wedge \cdots \wedge \vect{u}_n.
\]
\end{teiri}

これが行列式の「正体」である。第II部の深掘りで、行列式は
多重線形・交代・正規化の三性質で一意に特徴づけられると述べた。
外積代数は、その三性質を\textbf{代数構造として作り付けた空間}であり、
だから行列式が定義から自動的に流れ出てくる。
$\det(fg) = \det f \cdot \det g$ も、「合成変換の誘導する定数倍は
定数倍の積」という一行で証明が終わる。基底も成分も置換も使わない。
第I部で $ad - bc$ という謎の式として出会った行列式の、
これが最も高い場所から見た姿である。

\begin{itembox}[l]{深掘り:面積素・体積素 ― 微分形式への入口}
$\vect{u} \wedge \vect{v}$ の係数が「$\vect{u}, \vect{v}$ の張る
平行四辺形の符号付き面積」であることは、上の計算($ad - bc$)が示している。
$\R^3$ では $\Lambda^2 \R^3$ が $\binom{3}{2} = 3$ 次元になり、
$\vect{u} \wedge \vect{v}$ の $3$ つの係数は、高校で学んだ
\textbf{外積(クロス積)$\vect{u} \times \vect{v}$ の成分と一致する}
(クロス積が $3$ 次元でしか定義できなかった理由は、
$\binom{n}{2} = n$ となる $n$ が $3$ だけだからである)。

ウェッジ積を微分と組み合わせた\textbf{微分形式}の理論では、
$dx \wedge dy$ が「面積素」、$dx \wedge dy \wedge dz$ が「体積素」を表し、
重積分の変数変換でヤコビアン(行列式)が現れる理由、
ガウスの発散定理とストークスの定理が一本の定理
$\int_{\partial M} \omega = \int_M d\omega$ に統一される仕組みが、
すべて外積代数の上で展開される。線形代数の終点は、
多変数の微積分学・幾何学の出発点なのである。
\end{itembox}

\begin{mondai}\label{q:wedge}
　(1) $\vect{u} = \mat{1 \\ 2}$, $\vect{v} = \mat{3 \\ 4}$ について
$\vect{u} \wedge \vect{v}$ を $\vect{e}_1 \wedge \vect{e}_2$ の定数倍として
求め、係数が $\det\mat{1 & 3 \\ 2 & 4}$ に一致することを確かめよ。\\
　(2) $\R^3$ で $\vect{u} = \vect{e}_1 + \vect{e}_2$,
$\vect{v} = \vect{e}_2 + \vect{e}_3$ の外積 $\vect{u} \wedge \vect{v}$ を、
基底 $\vect{e}_1 \wedge \vect{e}_2$, $\vect{e}_1 \wedge \vect{e}_3$,
$\vect{e}_2 \wedge \vect{e}_3$ で展開せよ。
\end{mondai}
